\section{Limitations}
\label{sec:limitations}

Our study is deliberately controlled and correspondingly narrow. All token-based models use a frozen DINOv3 ViT-B/16 whose features are compressed to 256 dimensions and 8 bits; GOLF's hidden-test result relies on LoRA-adapted ViT-H+ features at higher resolution and on region sampling, which we do not reproduce, so absolute errors are not comparable to the leaderboard and our claims concern differences at a fixed representation. The evaluation server is closed and the test labels are withheld, so all numbers are on held-out training subjects; the official baseline is compared through the organisers' published number on the identical split rather than a retrained model. Two folds and two seeds for the main comparison give a range, not a variance, and the ablations and the second fold are single runs; differences below about half a millimetre should not be interpreted. The geometric head's predicted scale saturated on fold A and we did not intervene, because adjusting the clamp after inspecting held-out results would have compromised the protocol; a wider clamp, a warm start of the scale head or a separate learning rate are untested remedies. The geometric head presupposes a known rigid mesh and a given category; articulated or unknown objects would require a learned surface, and a category classifier would only partially close the gap. Our visibility labels ignore occlusion by the other hand. Ground-truth hand landmarks and object poses are themselves estimates with median errors of 6 to 8 and 1.5 to 3.5\,mm~\cite{rim2026show3d}, which bounds what can be measured at contact. HOT3D-eval has 538 frames from three participants. Temporal context is available at test time and unused. Every experiment ran on one 16\,GB laptop, which set the scale of the study; the same code runs unchanged on larger hardware.
